\section{Protected integration and orientation character attack--heal}
\label{sec:protected-integration-orientation-attack-heal}

\providecommand{\Z}{\mathbb Z}
\providecommand{\C}{\mathbb C}
\providecommand{\Hilb}{\operatorname{Hilb}}
\providecommand{\Ext}{\operatorname{Ext}}
\providecommand{\RHom}{\operatorname{RHom}}
\providecommand{\Perf}{\operatorname{Perf}}
\providecommand{\Crit}{\operatorname{Crit}}
\providecommand{\Prim}{\operatorname{Prim}}
\providecommand{\sdim}{\operatorname{sdim}}
\providecommand{\BPS}{\mathrm{BPS}}
\providecommand{\DT}{\mathrm{DT}}
\providecommand{\PT}{\mathrm{PT}}
\providecommand{\OP}{\mathrm{OP}}
\providecommand{\nuDelta}{\nu_{\Delta_5}}
\providecommand{\Gammaeff}{\Gamma_{\mathrm{eff}}}
\providecommand{\Hall}{\mathcal H}

\subsection{Claim attacked}

The scalar theorem gives
\[
  Z_{K3\times E}^{\OP}
  =
  Z_{K3\times E}^{\DT}
  =
  -4096\,\Delta_5^{-2},
\]
and the denominator normalization gives
\[
  \operatorname{den}(\mathfrak g_{\Delta_5})
  =
  64^{-1}\Delta_5(2Z).
\]
The attacked upgrade is stronger:

\begin{quote}
There is a protected Hall or cohomological integration map for
\(K3\times E\), with an orientation datum whose automorphic/reflection
character is exactly \(\nuDelta\).
\end{quote}

This upgrade is not a consequence of the scalar OP/DT theorem.  The
correct healing is to separate three objects:

\begin{enumerate}
\item the theorem-level scalar reduced integration map;
\item the theorem-level automorphic character of the Borcherds
determinant \(\Delta_5\);
\item the open protected-orientation datum which would identify the two
inside a compact Hall or cohomological Hall realization.
\end{enumerate}

\subsection{Source and local anchors}

\begin{center}
\small
\begin{tabular}{p{0.31\textwidth}p{0.60\textwidth}}
\hline
Anchor & Load-bearing content \\
\hline
\texttt{main.tex:547--577} &
\(\Delta_5\) has weight \(5\), \(\Delta_5^2=\Delta_{10}\), and
\(\Delta_5|_5g=\nuDelta(g)\Delta_5\). \\
\texttt{main.tex:581--585} &
The theta-normalization is
\([q^{1/2}r^{1/2}s^{1/2}]\Delta_5=64\); all coefficients lie in
\(64\Z\). \\
\texttt{main.tex:613--625} &
\(\Delta_5\) has the diagonal divisor with multiplicity one and is
anti-invariant under the three diagonal real-root reflections:
\(\Delta_5(s_{\delta_i}Z)=-\Delta_5(Z)\). \\
\texttt{main.tex:1228--1350} &
On \(W^{(2)}(\Lambda^{2,1}_{II})\),
\[
  \Delta_5(w\cdot Z)=\det(w)\Delta_5(Z),
\]
and the simple type-II real-root reflections have
\(\nuDelta(s_{\delta_i})=-1\). \\
\texttt{main.tex:1834--1954} &
The compact Hall construction is conditional on a compact Hall category,
HN finiteness, vanishing-cycle coefficients, Thom--Sebastiani
orientation transport, protected integration, and an orientation
character \(\epsilon_o\) satisfying
\(\epsilon_o|_{W^{(2)}}=\nuDelta|_{W^{(2)}}\). \\
\texttt{main.tex:2059--2095} &
The determinant/scalar square fixes signed exponents, not compact Hall
brackets, compact extension correspondences, or orientation constants. \\
\texttt{/tmp/igusa-cell8-sources/bryan/K3xE.tex:264--285} &
Reduced DT invariants are Behrend-weighted Euler characteristics of
\(\Hilb(X)/E\). \\
\texttt{/tmp/igusa-cell8-sources/bryan/K3xE.tex:333--364} &
The Bryan generating series uses
\(\widetilde q^{h-1}q^{d-1}(-p)^n\) and conjecturally gives
\(-1/\chi_{10}\). \\
\texttt{/tmp/igusa-cell8-sources/obred/reducedK3xE.tex:288--322} &
Reduced stable-pair invariants are Behrend-weighted orbifold Euler
characteristics of \(P_n(X,(\beta,d))/E\), equal to the reduced
virtual-class incidence invariants. \\
\texttt{/tmp/igusa-cell8-sources/obred/reducedK3xE.tex:432--466} &
Reduced DT invariants are defined by Behrend-weighted Euler
characteristics of Hilbert quotients, and
\(\DT_{(\beta,d)}(q)=\PT_{(\beta,d)}(q)\) for nonzero \(\beta\). \\
\texttt{/tmp/igusa-cell8-sources/op/main.tex:3212--3260} &
The OP/PT branch uses the reduced PT Behrend-weighted quotient
normalization and the GW/PT change \(y=-e^{iu}\). \\
\texttt{/tmp/igusa-cell8-sources/bridgeland/hall.tex:1588--1743} &
Joyce--Bridgeland semi-classical Hall integration maps weighted Euler
characteristics to a Poisson torus; Behrend weighting gives the DT sign
\(\sigma=-1\).  This is general Hall technology, not a constructed
compact \(K3\times E\) Hall atlas here. \\
\texttt{agent\_material/05\_normalization\_chain.tex} &
\(D_5=64^{-1}\Delta_5\),
\(\chi_{10}^{\OP}=D_5^2=4096^{-1}\Delta_5^2\), hence
\(Z^{\OP}=-4096\Delta_5^{-2}\). \\
\texttt{agent\_material/10\_compact\_hall\_atlas\_attack\_heal.tex} &
The compact Hall atlas, strong orientation, and character comparison
\(\epsilon_o=\nuDelta\) are still construction data, not established
consequences of the reduced quotient scalar theory. \\
\hline
\end{tabular}
\end{center}

\subsection{The integration map that actually exists}

\begin{proposition}[Reduced scalar integration]
\label{prop:reduced-scalar-integration-only}
For \(X=S\times E\), \(S\) a K3 surface, the OP/DT/PT theorem uses a
coefficientwise scalar integration map
\[
  [M]\longmapsto \int_M \nu_M\,de,
\]
applied to the reduced quotient spaces or quotient stacks
\[
  \Hilb^n(X,(\beta,d))/E,
  \qquad
  P_n(X,(\beta,d))/E.
\]
It lands in integers coefficient by coefficient.  It is the
Behrend-weighted, reduced, quotient Euler characteristic; in the
stable-pair branch it agrees with the reduced virtual-class incidence
definition of Oberdieck.
\end{proposition}

\begin{proof}
Bryan defines the reduced Hilbert-scheme DT invariant by
\[
  \DT_{h,d,n}^{\mathrm{red}}
  =
  \int_{\Hilb^{\beta_h+dE,n}(X)/E}\nu\,de,
\]
because the unreduced invariant on \(\Hilb(X)\) vanishes under the free
translation action of \(E\).  Oberdieck defines the stable-pair quotient
invariant by the Behrend-weighted orbifold Euler characteristic of
\(P_n(X,(\beta,d))/E\), proves that it equals the reduced incidence
invariant, and proves the reduced DT/PT equality
\[
  \DT_{(\beta,d)}(q)=\PT_{(\beta,d)}(q)
\]
for nonzero \(\beta\).  These statements are coefficientwise scalar
identities.  They do not require, and do not construct, a primitive
Hall algebra object.
\end{proof}

\begin{corollary}[What is not supplied]
\label{cor:what-is-not-supplied}
The scalar theorem does not supply any of the following objects:
\[
\begin{gathered}
  \Hall_{K3\times E}^{\mathrm{comp}},\qquad
  \text{compact extension correspondences},\qquad
  \text{HN-finite sector truncations},\\
  \text{a critical CoHA with vanishing-cycle coefficients},\qquad
  \text{a protected primitive projection }\Pi_{\Prim},\\
  \text{an orientation character }\epsilon_o
  \text{ with }\epsilon_o=\nuDelta .
\end{gathered}
\]
Thus the theorem-level conclusion is scalar:
\[
  Z_{K3\times E}^{\DT}=Z_{K3\times E}^{\OP}
  =-4096\Delta_5^{-2}.
\]
It is not a theorem-level construction of
\(\mathfrak g_{\Delta_5}\) from compact protected BPS primitives.
\end{corollary}

\begin{proof}
Joyce--Bridgeland Hall integration gives, in a suitable Calabi--Yau
threefold Hall setting, a semi-classical map from Hall elements to a
Poisson torus by weighted Euler characteristic, with \(\sigma=-1\) for
Behrend-weighted DT integration.  The local source itself treats the
motivic vanishing-cycle enhancement as additional technology, and its
simplified hypotheses do not by themselves give the reduced
\(E\)-quotient compact atlas needed here.  The manuscript also records
the missing inputs explicitly: compact extension correspondences,
Thom--Sebastiani orientation transport, protected integration, and
primitive extraction.  Hence the scalar quotient integration cannot be
rebranded as protected Hall integration.
\end{proof}

\subsection{Orientation character: theorem, criterion, open part}

\begin{proposition}[Automorphic character is proved]
\label{prop:automorphic-character-proved}
The character \(\nuDelta\) is theorem-level automorphic data of the
normalized Borcherds/Maass form \(\Delta_5\).  On the type-II Weyl
subgroup recorded in the manuscript,
\[
  \nuDelta(s_{\delta_i})=-1
  \qquad (i=1,2,3),
\]
and
\[
  \Delta_5(w\cdot Z)=\det(w)\Delta_5(Z)
  \qquad
  \bigl(w\in W^{(2)}(\Lambda^{2,1}_{II})\bigr).
\]
The scalar square \(\Delta_5^2=\Delta_{10}\) forgets this sign.
\end{proposition}

\begin{proof}
This is the Maass/Borcherds character statement recorded in
\texttt{main.tex:547--577} and the reflection computation recorded in
\texttt{main.tex:1228--1350}.  The simple diagonal reflections are
anti-invariant for \(\Delta_5\), while the square is scalar.  No
enumerative orientation input is used in this proof.
\end{proof}

\begin{definition}[Protected orientation matching datum]
\label{def:protected-orientation-matching-datum}
A compact protected integration/orientation datum realizing the
\(\Delta_5\) denominator would consist of:

\begin{enumerate}
\item an extension-closed compact Hall or critical CoHA object
\(\Hall_{K3\times E}^{\mathrm{comp}}\) in a strict stability sector,
with HN-finite truncations and support property;
\item derived critical charts
\[
  \mathfrak M_\gamma|_U\simeq[\Crit(f_{\gamma,U})/G_{\gamma,U}]
\]
for the relevant charge stacks;
\item a strong orientation datum \(o\), namely a square root of the
determinant line
\[
  o_\gamma^{\otimes 2}
  \simeq
  \operatorname{sdet}\RHom(\mathcal E_\gamma,\mathcal E_\gamma)
\]
with exact-triangle and Thom--Sebastiani coherences;
\item equivariant descent through the reduced \(E\)-quotient;
\item a protected integration map
\[
  I^{\mathrm{prot}}_{\sigma,o}:
  \Hall_{K3\times E}^{\mathrm{comp}}
  \longrightarrow
  \widehat{\C[\Gammaeff]}
\]
or its semi-classical/protected primitive shadow;
\item a primitive projection satisfying
\[
  \sdim \Prim_\gamma=f(nm,l)
  \qquad
  \text{for }\gamma=(n,l,m);
\]
\item a symmetry action on the oriented theory whose induced
orientation character
\[
  \epsilon_o:W^{(2)}(\Lambda^{2,1}_{II})\longrightarrow\{\pm1\}
\]
satisfies the comparison
\[
  \epsilon_o=\nuDelta
  \quad\text{on }W^{(2)}(\Lambda^{2,1}_{II}).
\]
\end{enumerate}
\end{definition}

\begin{proposition}[Orientation match is not a normalization choice]
\label{prop:not-normalization-choice}
The equality \(\epsilon_o=\nuDelta\) is not obtained merely by choosing
the leading sign of \(\Delta_5\).  Multiplying \(\Delta_5\) by a nonzero
constant changes neither its automorphic character nor its reflection
anti-invariance.  The normalization
\[
  [q^{1/2}r^{1/2}s^{1/2}]\Delta_5=64
\]
chooses the theta-product section.  It does not construct a strong
orientation on the Hall side and does not prove that the orientation
monodromy is \(\nuDelta\).
\end{proposition}

\begin{proof}
Characters are ratios:
\[
  \frac{(c\Delta_5)|_5g}{c\Delta_5}
  =
  \frac{\Delta_5|_5g}{\Delta_5}
  =
  \nuDelta(g).
\]
Thus the scalar constant \(c\) cancels.  A geometric orientation
character instead comes from transporting a chosen square root of a
determinant line around symmetries or wall reflections.  That
transport has not been built for the compact \(K3\times E\) Hall object.
\end{proof}

\subsection{How the diagonal reflection sign should arise}

The automorphic fact is exact:
\[
  \operatorname{div}(\Delta_5)
  =
  \operatorname{Sp}_4(\Z)\cdot H_{\mathrm{diag}}
  \quad\text{with multiplicity }1,
\]
and each simple diagonal real-root reflection \(s_{\delta_i}\) satisfies
\[
  \Delta_5(s_{\delta_i}Z)=-\Delta_5(Z).
\]

The expected orientation mechanism is the following.  The diagonal
Humbert divisor is the wall where a single real-root normal mode becomes
massless.  A square root of the determinant of the self-dual deformation
complex behaves like a Pfaffian line.  Reflection in the normal direction
should reverse that oriented one-dimensional normal factor, giving a
sign \(-1\) on the square root and no sign on its square.  This is the
geometric shape of
\[
  \epsilon_o(s_{\delta_i})=-1,
  \qquad
  \epsilon_o^2(s_{\delta_i})=1.
\]

This paragraph is a criterion, not a proof.  A proof would have to
construct the oriented critical Hall atlas, identify the diagonal
wall-crossing normal mode with the simple type-II real root
\(\delta_i\), and compute the induced monodromy of the orientation line.
The current paper proves the automorphic sign of \(\Delta_5\); it does
not prove this Pfaffian-orientation computation from \(K3\times E\)
geometry.

\subsection{Attack on the constants}

\begin{center}
\small
\begin{tabular}{p{0.20\textwidth}p{0.69\textwidth}}
\hline
Constant/sign & Status \\
\hline
\(64\) &
Theta-product normalization:
\([q^{1/2}r^{1/2}s^{1/2}]\Delta_5=64\).  The monic denominator
section is
\[
  D_5=64^{-1}\Delta_5.
\]
This constant is not a Hall integration constant, not an orientation
sign, and not produced by the Weyl vector. \\
\(-4096\) &
\[
  \chi_{10}^{\OP}=D_5^2=4096^{-1}\Delta_5^2,
  \qquad
  Z^{\OP}=-(\chi_{10}^{\OP})^{-1}.
\]
Thus
\[
  Z^{\OP}=-4096\Delta_5^{-2}.
\]
Oberdieck's DT/PT theorem transfers this same constant to reduced DT.
The minus sign is the OP curve-counting convention, not the reflection
character of \(\Delta_5\). \\
\(64^{-1}\Delta_5(2Z)\) &
The denominator uses the monic section \(D_5\) and the doubled
exponential convention for the root lattice:
\[
  \operatorname{den}(\mathfrak g_{\Delta_5})
  =
  D_5(2Z)
  =
  64^{-1}\Delta_5(2Z).
\]
The factor \(64^{-1}\) removes the theta leading coefficient; the
argument \(2Z\) is the root-coordinate normalization. \\
Weyl vector &
The Weyl vector fixes the leading monomial
\[
  q^{1/2}r^{1/2}s^{1/2}
\]
or its doubled-coordinate version.  It does not create the numerical
coefficient \(64\), and it does not explain the OP sign. \\
\hline
\end{tabular}
\end{center}

\subsection{Exact paper language}

The theorem-level statement should be:

\begin{quote}
For \(X=K3\times E\), the reduced DT/PT/OP scalar integration is the
Behrend-weighted reduced quotient Euler characteristic on
\(\Hilb(X)/E\) and \(P_n(X,\beta)/E\).  With the theta normalization
\([q^{1/2}r^{1/2}s^{1/2}]\Delta_5=64\), the primitive scalar partition
function is
\[
  Z_{X}^{\DT}=Z_X^{\OP}=-4096\Delta_5^{-2}.
\]
\end{quote}

The automorphic theorem should be:

\begin{quote}
The Borcherds determinant \(\Delta_5=\operatorname{Borch}(\phi_{0,1})\)
has Maass character \(\nuDelta\).  Its divisor is the diagonal Humbert
divisor with multiplicity one, and on the type-II Weyl subgroup
\[
  \Delta_5(w\cdot Z)=\det(w)\Delta_5(Z).
\]
In particular \(\nuDelta(s_{\delta_i})=-1\) for the simple diagonal
real-root reflections.
\end{quote}

The protected realization should not be stated as a theorem.  It should
be stated as the following criterion/open problem:

\begin{quote}
A compact protected \(K3\times E\) realization of
\(\mathfrak g_{\Delta_5}\) is an oriented critical Hall or CoHA object
with reduced \(E\)-quotient descent, protected integration
\(I^{\mathrm{prot}}_{\sigma,o}\), primitive extraction, and a strong
orientation datum \(o\) whose induced reflection character
\(\epsilon_o\) satisfies
\[
  \epsilon_o|_{W^{(2)}(\Lambda^{2,1}_{II})}
  =
  \nuDelta|_{W^{(2)}(\Lambda^{2,1}_{II})}.
\]
If such a datum is constructed and if
\(\sdim\Prim_{(n,l,m)}=f(nm,l)\), then its denominator is
\[
  \operatorname{den}(\mathfrak g_{\Delta_5})
  =
  64^{-1}\Delta_5(2Z).
\]
The existence of this oriented compact protected integration datum is
not proved by the scalar OP/DT theorem and remains a realization
problem.
\end{quote}

\subsection{Claim-status recommendation}

\begin{center}
\small
\begin{tabular}{p{0.37\textwidth}p{0.13\textwidth}p{0.39\textwidth}}
\hline
Claim & Status & Reason \\
\hline
Reduced OP/DT scalar square
\(-4096\Delta_5^{-2}\) &
Theorem &
Oberdieck DT/PT plus OP/PT/GW comparison and OP
\(-1/\chi_{10}\), with the manuscript's \(64\)-normalization. \\
Behrend-weighted reduced quotient integration exists &
Theorem &
Defined on \(\Hilb(X)/E\) and \(P_n(X,\beta)/E\); stable-pair quotient
equals reduced virtual-class incidence. \\
Motivic/CoHA/perverse-sheaf protected integration for this compact
realization &
Open datum &
General Hall and CoHA technologies exist, but the compact
\(K3\times E\) Hall atlas, quotient descent, primitive projection, and
orientation transport are not constructed here. \\
\(\nuDelta\) is the automorphic/reflection character of \(\Delta_5\) &
Theorem &
Maass/Borcherds character and reflection computation in the manuscript. \\
There is a geometric orientation datum with
\(\epsilon_o=\nuDelta\) &
Open problem /
criterion &
Requires a strong orientation and a monodromy computation on the
oriented compact Hall theory.  Scalar normalization cannot supply it. \\
The diagonal reflection sign comes from orientation &
Expected mechanism &
The Pfaffian-line explanation is the right local model, but it has not
been derived from a constructed \(K3\times E\) orientation datum. \\
Constants \(64\), \(4096\), and \(-\) are Hall-orientation constants &
False &
\(64\) is theta normalization, \(4096=64^2\), and the minus sign is the
OP scalar convention. \\
\hline
\end{tabular}
\end{center}

\subsection{Files changed, checks run, remaining questions}

Files changed:
\[
  \texttt{agent\_material/12\_protected\_integration\_orientation\_attack\_heal.tex}.
\]

Checks run:
\begin{enumerate}
\item Loaded the required repository discipline files:
\(\texttt{CLAUDE.md}\), \(\texttt{AGENTS.md}\), and
\(\texttt{\string~/ecosystem/INVARIANTS.md}\).
\item Read \(\texttt{main.tex}\), \(\texttt{proj.bib}\), and the local
OP/DT/Hall agent-material notes named in the mission scope.
\item Checked the local primary-source TeX for Bryan, Oberdieck,
Oberdieck--Pandharipande/Oberdieck--Pixton, and Joyce--Bridgeland Hall
integration anchors listed above.
\item Ran an \(\texttt{rg}\) sanity scan for overclaiming and attribution
trigger phrases in this report; no hits.
\item No full paper build was run: this report is not included by
\(\texttt{main.tex}\), and no main manuscript file was edited.
\end{enumerate}

Remaining open questions:
\begin{enumerate}
\item Construct an extension-closed compact Hall or critical CoHA object
for the reduced \(K3\times E\) sector, with HN-finite sector
truncations.
\item Prove \(E\)-equivariant descent of the strong orientation and the
vanishing-cycle coefficient system through the reduced quotient.
\item Construct the protected primitive projection and prove
\[
  \sdim\Prim_{(n,l,m)}=f(nm,l).
\]
\item Compute the orientation-line monodromy under the simple diagonal
reflections and prove
\[
  \epsilon_o(s_{\delta_i})=-1.
\]
\item Extend that computation from the simple reflections to the full
comparison
\[
  \epsilon_o|_{W^{(2)}(\Lambda^{2,1}_{II})}
  =
  \nuDelta|_{W^{(2)}(\Lambda^{2,1}_{II})}.
\]
\end{enumerate}
